La capacidad nominal por sí sola no determina el vaciado. Un servicio a plena capacidad mientras haya trabajo vacía el estado continuo en \(3\) min y lo mantiene en cero. Otro modelo, con salida \(\min(q/\tau,10)\) y \(\tau=1\) min, conserva trabajo pendiente y tiende a \(6\) piezas. La conclusión debe acompañarse de la ley elegida y su interpretación.
